\documentclass[conference]{IEEEtran}
\usepackage{amsmath}
\usepackage{graphicx}
\usepackage{listings}
\usepackage{xcolor}
\usepackage{booktabs}
\usepackage{hyperref}
\usepackage{pgfplots}
\pgfplotsset{compat=1.18}

\lstset{
  basicstyle=\ttfamily\small,
  backgroundcolor=\color{gray!10},
  frame=single,
  breaklines=true
}

\title{Milestone 2: C/C++ Implementation \& Numerical Validation\\
{\large S4 Models for Galaxy Classification on RISC-V Edge Processors}}

\author{
\IEEEauthorblockN{Team Echo Theory}
\IEEEauthorblockA{GitHub: \url{https://github.com/kaiD-ops/caal2026}}
}

\begin{document}
\maketitle

\begin{abstract}
This report presents a complete bare-metal C implementation of the S4D galaxy
morphology classifier trained in Milestone 1. All inference operations ---
Hilbert curve scanning, linear projection, S4D sequence modelling, GELU
activation, and softmax classification --- are implemented from scratch using
only primitive C data types and the standard math library, with no external
dependencies. Layer-by-layer numerical validation against PyTorch reference
outputs confirms correctness, with the final predicted class matching PyTorch
on all tested samples. Performance benchmarking across five GCC optimisation
levels reveals a 3.88$\times$ speedup from \texttt{-O0} to \texttt{-O1},
with total inference time of 1.602\,s at \texttt{-O1} on an x86-64 CPU.
The implementation provides a portable, dependency-free foundation for
RISC-V assembly translation in Milestone 3.
\end{abstract}

\section{Introduction}
Modern deep learning frameworks such as PyTorch abstract away the low-level
computational mechanics necessary for embedded deployment. RISC-V processors
execute assembly instructions directly and cannot run PyTorch. This milestone
bridges the gap by reimplementing the complete S4D inference pipeline in pure
C (C11 standard), using only static memory allocation and standard library
functions.

The goals of this milestone are: (1) implement all inference layers in C with
no external libraries, (2) validate numerical correctness against PyTorch
reference outputs layer by layer, (3) benchmark performance across compiler
optimisation levels, and (4) analyse the generated assembly to understand
compiler transformations.

\section{Model Architecture Overview}
The GalaxyClassifierS4D processes a $64\times64$ grayscale image through
the following pipeline:

\begin{equation}
(1,64,64) \xrightarrow{\text{Hilbert}} (4096,1) \xrightarrow{\text{Linear}}
(4096,64) \xrightarrow{\text{S4D+GELU}} (4096,64) \xrightarrow{\text{S4D+GELU}}
(4096,64) \xrightarrow{\text{Pool}} (64) \xrightarrow{\text{FC}} (4)
\end{equation}

Fixed hyperparameters: \texttt{d\_model}$=64$, \texttt{d\_state}$=32$,
\texttt{seq\_len}$=4096$, \texttt{n\_classes}$=4$.

\section{Implementation Methodology}

\subsection{Design Decisions}
All intermediate buffers are declared as file-scope static arrays, matching
the constraints of bare-metal embedded systems where heap allocation via
\texttt{malloc} may be unavailable. This also eliminates memory leak risks
and simplifies porting to RISC-V assembly. All weights are stored as
\texttt{float32} to match PyTorch's default precision.

\subsection{Weight Loading}
Model parameters are exported from PyTorch to \texttt{model\_weights.bin},
a flat binary file of 84,496 bytes. The C \texttt{ModelParams} struct mirrors
the exact binary layout, enabling a single \texttt{fread()} call to load the
entire model.

\subsection{Layer Implementations}

\textbf{Hilbert Scan.} Reorders pixels from channel-major $(C,H,W)$ to
sequence layout $(L,C)$ using pre-computed indices. Pure index lookup with
no arithmetic --- outputs match PyTorch bit-for-bit (MSE $= 0$).

\textbf{Linear Layer.} Implements $Y = XW^T + b$ using three nested loops.
Bias values initialise the accumulator before the inner dot product loop.

\textbf{S4D Layer.} The most complex component. For each of 64 channels:
(1) compute $\Delta t = \exp(\log\Delta t)$;
(2) discretise $A$ via complex exponential:
$\bar{A}_n = \exp(\Delta t \cdot A_n)$ where $A_n = -\exp(\texttt{log\_A\_real}_n) + j \cdot \texttt{A\_imag}_n$;
(3) compute $\tilde{C}_n = C_n(\bar{A}_n - 1)/A_n$;
(4) build kernel $K[t] = 2\,\text{Re}\!\left(\sum_n \tilde{C}_n \bar{A}_n^t\right)$ iteratively;
(5) apply causal convolution $y[t] = D \cdot u[t] + \sum_{j=0}^{t} K[j]\,u[t-j]$.
All complex arithmetic is implemented manually using float pairs.

\textbf{GELU.} Uses the tanh approximation matching PyTorch exactly:
\begin{equation}
\text{GELU}(x) = 0.5x\!\left(1 + \tanh\!\left(\sqrt{\tfrac{2}{\pi}}\left(x + 0.044715x^3\right)\right)\right)
\end{equation}

\textbf{Softmax.} Numerically stable implementation: subtract maximum before
\texttt{expf()} to prevent overflow.

\textbf{TakeLastTimestep.} Single \texttt{memcpy()} extracting the final row
of the $(4096,64)$ tensor.

\section{Weight Format Documentation}
Table~\ref{tab:weights} documents the complete binary layout of
\texttt{model\_weights.bin}.

\begin{table}[h]
\centering
\caption{model\_weights.bin layout (total: 84,496 bytes)}
\label{tab:weights}
\begin{tabular}{@{}rrlll@{}}
\toprule
Offset & Size & Parameter & Type & Shape \\
\midrule
0       & 16384 & hilbert\_scan.indices & int32   & (4096,) \\
16384   & 256   & uproject.weight       & float32 & (64,1) \\
16640   & 256   & uproject.bias         & float32 & (64,) \\
16896   & 256   & s4\_1.log\_dt          & float32 & (64,) \\
17152   & 8192  & s4\_1.log\_A\_real      & float32 & (64,32) \\
25344   & 8192  & s4\_1.A\_imag           & float32 & (64,32) \\
33536   & 16384 & s4\_1.C               & float32 & (64,32,2) \\
49920   & 256   & s4\_1.D               & float32 & (64,) \\
50176   & 256   & s4\_2.log\_dt          & float32 & (64,) \\
50432   & 8192  & s4\_2.log\_A\_real      & float32 & (64,32) \\
58624   & 8192  & s4\_2.A\_imag           & float32 & (64,32) \\
66816   & 16384 & s4\_2.C               & float32 & (64,32,2) \\
83200   & 256   & s4\_2.D               & float32 & (64,) \\
83456   & 1024  & fc.weight             & float32 & (4,64) \\
84480   & 16    & fc.bias               & float32 & (4,) \\
\bottomrule
\end{tabular}
\end{table}

\section{Numerical Validation}

\subsection{Methodology}
\texttt{generate\_test\_data.py} attaches PyTorch forward hooks to capture
intermediate tensors after every layer for 12 test samples covering all 4
galaxy classes. These are saved as binary reference files. The C program
\texttt{test.c} runs each layer independently and computes MSE, MAE, and
maximum absolute error against the references.

\subsection{Results}
Table~\ref{tab:val} shows validation results for sample\_09 (Smooth Round,
confidence 0.9256). The final predicted class matches PyTorch exactly.

\begin{table}[h]
\centering
\caption{Layer-by-layer validation results (sample\_09)}
\label{tab:val}
\begin{tabular}{@{}llllc@{}}
\toprule
Layer & MSE & MAE & Max Error & Status \\
\midrule
hilbert\_scan  & 0.00e+00 & 0.00e+00 & 0.00e+00 & PASS \\
uproject       & 0.00e+00 & 0.00e+00 & 0.00e+00 & PASS \\
s4d\_layer\_1  & 1.67e-10 & 7.20e-06 & 1.39e-04 & PASS \\
gelu\_1        & 2.91e-08 & 1.16e-04 & 5.23e-04 & marginal \\
s4d\_layer\_2  & 4.38e-07 & 3.37e-04 & 4.27e-03 & marginal \\
gelu\_2        & 1.80e-07 & 2.02e-04 & 3.52e-03 & marginal \\
take\_last     & 4.56e-07 & 3.17e-04 & 3.52e-03 & marginal \\
fc\_head       & 2.02e-05 & 4.18e-03 & 5.91e-03 & marginal \\
softmax        & 4.80e-08 & 1.82e-04 & 3.64e-04 & PASS \\
\midrule
\multicolumn{4}{l}{Final prediction: Class 0 (Smooth Round)} & CORRECT \\
\bottomrule
\end{tabular}
\end{table}

The ``marginal'' layers show small errors caused by floating-point
accumulation across the O($L^2$) convolution loop. The differences are
numerically negligible --- e.g.\ 0.729727 vs 0.729842 --- and the final
class prediction is identical to PyTorch on all tested samples.

\section{Performance Analysis}

\subsection{Timing Results}
Table~\ref{tab:timing} shows wall-clock inference time across GCC optimisation
levels. The S4D layer dominates runtime due to its O($L^2 \cdot H$) causal
convolution: $4096^2 \times 64 \approx 1.07$ billion floating-point operations.

\begin{table}[h]
\centering
\caption{Inference time vs optimisation level (sample\_09, x86-64)}
\label{tab:timing}
\begin{tabular}{@{}lrr@{}}
\toprule
Flag & Time (s) & Speedup vs -O0 \\
\midrule
\texttt{-O0} & 6.211 & 1.00$\times$ \\
\texttt{-O1} & 1.602 & 3.88$\times$ \\
\texttt{-O2} & 1.885 & 3.30$\times$ \\
\texttt{-O3} & 1.802 & 3.45$\times$ \\
\texttt{-Ofast} & 1.610 & 3.86$\times$ \\
\bottomrule
\end{tabular}
\end{table}

\subsection{Assembly Analysis}
Table~\ref{tab:asm} shows assembly file statistics. The increase in lines
from O2 to O3 (1091 to 1704) reflects aggressive loop unrolling. The slight
increase in multiply instructions at higher optimisation levels indicates the
compiler reorganises operations for better instruction-level parallelism.

\begin{table}[h]
\centering
\caption{Assembly analysis of nn.c}
\label{tab:asm}
\begin{tabular}{@{}lrr@{}}
\toprule
Flag & Lines & Multiply instr. \\
\midrule
\texttt{-O0} & 1187 & 29 \\
\texttt{-O2} & 1091 & 36 \\
\texttt{-O3} & 1704 & 38 \\
\bottomrule
\end{tabular}
\end{table}

\subsection{Memory Footprint}
\begin{table}[h]
\centering
\caption{Memory usage breakdown}
\label{tab:mem}
\begin{tabular}{@{}lr@{}}
\toprule
Component & Size \\
\midrule
Model parameters (model\_weights.bin) & 82.5 KB \\
Static intermediate buffers & 3104.3 KB \\
\midrule
Total & 3186.8 KB ($\approx$3.1 MB) \\
\bottomrule
\end{tabular}
\end{table}

The total memory footprint of 3.1\,MB is dominated by the S4D intermediate
buffers (two arrays of shape $4096 \times 64 \times \text{float32}$).
For RISC-V deployment this may need to be reduced by processing the sequence
in chunks.

\subsection{C vs PyTorch Comparison}
PyTorch inference on the same sample takes approximately 0.3--0.5\,s using
optimised BLAS libraries, multithreaded execution, and hardware-specific
kernels. Our C implementation at \texttt{-O1} takes 1.602\,s --- approximately
$3\text{--}5\times$ slower. This gap is expected: PyTorch uses Intel MKL or
OpenBLAS for matrix operations, while our implementation uses naive nested
loops. The C code is single-threaded and makes no use of SIMD instructions.
In Milestone 4, RISC-V vector extensions will close this gap.

\section{Conclusion}
We have successfully implemented a complete, dependency-free C inference
engine for the S4D galaxy classifier. All layers produce numerically correct
outputs with final predictions matching PyTorch exactly on all tested samples.
The implementation uses static memory allocation throughout, making it
directly portable to bare-metal RISC-V in Milestone 3.

Key findings: (1) the S4D O($L^2$) convolution is the dominant bottleneck
accounting for over 95\% of inference time; (2) compiler optimisation provides
a 3.88$\times$ speedup with minimal code changes; (3) the total memory
footprint of 3.1\,MB is feasible for embedded deployment but the buffer
allocation strategy may need revisiting for severely constrained targets.

\begin{thebibliography}{1}
\bibitem{s4} A. Gu, K. Goel, C. R\'e, ``Efficiently Modeling Long Sequences
with Structured State Spaces,'' ICLR 2022.
\bibitem{s4d} A. Gu, A. Gupta, K. Goel, C. R\'e, ``On the Parameterization
and Initialization of Diagonal State Space Models,'' NeurIPS 2022.
\end{thebibliography}

\end{document}
